\scanpage{70}
% No authorial number is visible at the head of this sheet.
Supposons en plus

d) toute flèche inversible de $\uline{\Sigma}$ est une identité
(p.ex. $\uline{\Sigma}$ est une catégorie définie par un ens
\uline{ordonné}).

Alors on conclut aisément que l'ens. pré-%
ordonné $I$ est \uline{ordonné}, et que tte \unsure{flèche|flèches}
de $\uline{\Sigma}$ qui n'est pas une identité i.e. $u\in\Sigma^{\unsure{*|<}}(1)$,
son image dans $I(1)$ n'est pas une identité;
\deleted{\unsure{i.e.}} Si $\uline{\Sigma}$ est associé à un
\deleted{\unsure{ens. ord.}}\inserted{espace top.}
ordonné $\Sigma=\Sigma_{(0)}$, cela signifie que
$\uline{\Sigma}\mathrel{\relbar\joinrel\relbar}\uline{I}$ est \inserted{une appl.}
\uline{strictement} croissante.
% pré- has a written line-end dash.  The upper-right ending of flèche is
% reworked, with a raised curl and sloping stroke; a final s remains open.
% The crossed i.e. and ens. ord. have obscured punctuation/letter joins.
% espace top. stands immediately above the cancelled ens. ord. after un.
% Sigma=Sigma_(0) follows ordonné on the rising main baseline, without an
% independent insertion leader or preceding comma.  The first map's shaft
% has a short downward terminal stub; both object symbols are underlined.
% The curved leader for une appl. enters immediately before strictement.

Revenons alors à un \deleted{catég} ens. ordonné topologique
\[
\uline{\Sigma}=(\Sigma=\Sigma(0),\quad
\underset{\underset{\Sigma\unsure{\times}\Sigma}{\cap}}{\Sigma(1)}
=\delta(\Sigma_{(0)})\amalg\Sigma^{<}(1))
\]
% The squeezed live Sigma= lies just inside the opening parenthesis.
% The first (0) and following (1) have baseline placement; the inner (0)
% in delta's argument is lowered.  The cap and product remain vertical.
% The product sign has a looped left part and crossing right arms.
satisfaisant aux conditions a) b) c) et $1^\circ$) de la page (27)
(correspondant \inserted{moralement} \unsure{à|a} un déploiement d'une
stratification locale équisingulier), et tel que les comp. connexes
de $\Sigma$ \deleted{\unsure{sont|soient}} de dim finie.
Considérons alors \unsure{la|le} \deleted{\unsure{factorisation}}\inserted{morphisme}
canonique $\uline{\Sigma}\unsure{:}\mathrel{\relbar\joinrel\relbar}\uline{I}$ avec $I=\pi_0(\Sigma)$,
% moralement has its own leader after correspondant; the possible grave
% on the following a lies within that leader zone.  The g-descender in
% équisingulier is continuous with the letter, not a separate underline.
% morphisme is above factorisation, not attached to the cancelled sont.
% The article's final small lobe leaves la/le open.  The cancelled word's
% right end reaches the paper edge.  The second map's two small tilted
% points are kept as an uncertain colon; its shaft has a dense terminal
% upright.  The target I is reworked, without an assigned older glyph.
et supposons de plus que pour $i<j$ ($i,j\in I$)
\[
\Sigma_{ij}\unsure{\longrightarrow|-}\Sigma_j\quad\hbox{soit \uline{mono} i.e. plongement fermé}
\]
% The short terminal point of the mono-map is filled and compact; the
% arrow/shaft contour remains doubtful.  fermé meets the right edge,
% with an accent but no independent terminal full stop visible.
On est alors sous les conditions de p. 28,
données a) b) c)
$(\unsure{\Sigma|\Sigma_i},\Sigma_{ij},\ s_{ij}\ b_{ij}:$
\diagram{\begin{tikzcd}[column sep=15mm,row sep=-4mm,execute at end picture={
% Two short, nearly parallel, slightly divergent shafts.  Their source
% heights differ.  The upper endpoint is dense and pointed; the lower
% endpoint contains a small loop, kept open in the drawing.
\draw[line width=.55pt,line cap=round] ([yshift=2pt]s70source.east) -- ([xshift=-2pt]s70upper.west);
\fill ([xshift=-3pt,yshift=.2pt]s70upper.west) --
([xshift=-2pt,yshift=1.4pt]s70upper.west) --
([xshift=-.9pt,yshift=.55pt]s70upper.west) --
([xshift=-1.4pt,yshift=-.65pt]s70upper.west) --
([xshift=-2.5pt,yshift=-.3pt]s70upper.west) -- cycle;
\draw[line width=.55pt,line cap=round] ([yshift=-2pt]s70source.east) -- ([xshift=-2.7pt]s70lower.west);
\draw[line width=.5pt] ([xshift=-2pt]s70lower.west) ellipse[x radius=.7pt,y radius=.6pt];
}]
& |[alias=s70upper]|\Sigma_i\\
|[alias=s70source]|\Sigma_{ij}&\\
& |[alias=s70lower]|\Sigma_j
\end{tikzcd}}$),$
% The first tuple Sigma has a compact raised foot-tip; a separate i cannot
% be securely isolated.  No independent comma is assigned between s_ij and b_ij.
% Pale bottom-left foliation, another hand.
\margin{69}
